\chapter{Аналитическая часть}

\section{Назначение разрабатываемого компилятора}

Разрабатываемый компилятор предназначен для трансляции программ на языке B в ассемблерный код архитектуры x86-64 с последующим получением исполняемого файла под операционную систему Linux. Входом компилятора является один или несколько файлов с исходным текстом программы. Выходом может быть:

\begin{itemize}
  \item ассемблерный файл с расширением \texttt{.s};
  \item объектный файл с расширением \texttt{.o};
  \item исполняемый файл, полученный после компоновки.
\end{itemize}

Компилятор реализован на языке C. Для сборки промежуточного ассемблерного текста используется системный ассемблер, для компоновки --- системный компилятор C, вызываемый как драйвер линковщика. Такой подход позволяет сосредоточить реализацию на трансляции языка B, не разрабатывая собственный ассемблер и компоновщик.

\section{Краткая характеристика языка B}

Согласно руководству Ken Thompson ``Users' Reference to B'', язык B предназначен для рекурсивных, преимущественно ненумерических задач системного программирования \cite{bmanual}. Язык имеет небольшую грамматику, простую модель хранения и богатую систему выражений.

Основные свойства языка B:

\begin{itemize}
  \item программа состоит из последовательности внешних определений;
  \item внешнее определение может задавать глобальный объект, вектор или функцию;
  \item функции объявляются именем, списком параметров и оператором-телом;
  \item локальные автоматические объекты объявляются ключевым словом \texttt{auto};
  \item внешние имена, используемые внутри функции, объявляются ключевым словом \texttt{extrn};
  \item язык оперирует машинными словами и исторически не требует строгой системы типов;
  \item указательная арифметика выражается через операции \texttt{*}, \texttt{\&} и индексирование;
  \item комментарии записываются в форме \texttt{/* ... */};
  \item escape-последовательности в строковых и символьных литералах начинаются с символа \texttt{*}.
\end{itemize}

В реализации курсовой работы базовая модель B адаптирована к x86-64. Основной единицей хранения является машинное слово целевой архитектуры. Внешние определения, автоматические объекты, векторы, функции и выражения разбираются в формах, описанных в руководстве по языку B.

\section{Теоретические основы трансляции}

Процесс компиляции можно рассматривать как последовательное преобразование программы из одной формы представления в другую. На первом этапе исходный текст является обычной последовательностью символов, в которой еще не выделены смысловые единицы языка. Лексический анализ группирует символы в токены: ключевые слова, идентификаторы, литералы, операции и разделители. После этого синтаксический анализ проверяет, что последовательность токенов соответствует грамматике языка, и строит внутреннее представление программы. В данной работе таким представлением является абстрактное синтаксическое дерево, поскольку оно сохраняет смысловую структуру программы, но не содержит лишних синтаксических деталей вроде скобок и разделителей.

Абстрактное синтаксическое дерево удобно для компилятора тем, что разные внешние формы записи могут быть сведены к одинаковой внутренней операции. Например, индексирование вектора в языке B выражает адресную арифметику, поэтому выражение обращения к элементу может быть представлено как вычисление адреса и последующее разыменование. Аналогично составное присваивание сначала разбирается как специальная форма токена, а затем превращается в дерево, в котором отдельно представлены чтение старого значения, бинарная операция и запись результата. Такая нормализация уменьшает количество специальных случаев на этапе генерации кода.

Генерация машинно-зависимого кода требует учета соглашений целевой платформы. Для x86-64 под Linux важны размер машинного слова, способ размещения локальных переменных в стековом кадре, регистры передачи аргументов, регистр возврата результата и правила сохранения регистров при вызовах функций. Компилятор данной работы генерирует ассемблерный текст, а затем передает его системному ассемблеру и компоновщику. Благодаря этому этап трансляции остается достаточно наглядным: результат можно просмотреть как обычный текстовый файл, сопоставить инструкции с исходными конструкциями языка B и отдельно проверить корректность вызовов внешних функций, таких как \texttt{printf}.

\section{Лексические элементы}

Лексический анализатор должен выделять из входного потока токены, приведенные в таблице \ref{tbl:tokens}.

\begin{longtable}{|p{4.2cm}|p{11.8cm}|}
  \caption{Основные лексические элементы языка}
  \label{tbl:tokens} \\
  \hline
  \textbf{Группа} & \textbf{Элементы} \\
  \hline
  \endfirsthead
  \hline
  \textbf{Группа} & \textbf{Элементы} \\
  \hline
  \endhead
  Идентификаторы & Имена функций, переменных, векторов и меток. Состоят из букв, цифр и символа подчеркивания, не начинаются с цифры. \\
  \hline
  Ключевые слова B & \texttt{auto}, \texttt{extrn}, \texttt{if}, \texttt{else}, \texttt{while}, \texttt{switch}, \texttt{case}, \texttt{return}, \texttt{goto}. \\
  \hline
  Литералы & Десятичные и восьмеричные целые константы, символьные константы из одного или двух символов, строковые литералы. \\
  \hline
  Операторы & \texttt{=}, \texttt{=+}, \texttt{=-}, \texttt{=*}, \texttt{=/}, \texttt{=\%}, \texttt{=|}, \texttt{=\&}, \texttt{=<<}, \texttt{=>>}, \texttt{===}, \texttt{=!=}, арифметические, битовые, логические, сравнительные и тернарный оператор. \\
  \hline
  Разделители & \texttt{;}, \texttt{,}, \texttt{:}, \texttt{?}, круглые, квадратные и фигурные скобки. \\
  \hline
\end{longtable}

Поддерживаемые escape-последовательности языка B приведены в таблице \ref{tbl:escapes}.

\begin{longtable}{|p{4cm}|p{12cm}|}
  \caption{Escape-последовательности языка B}
  \label{tbl:escapes} \\
  \hline
  \textbf{Последовательность} & \textbf{Значение} \\
  \hline
  \texttt{*0} & Нулевой символ. \\
  \hline
  \texttt{*e} & Признак конца файла; в реализации представлен значением \texttt{EOF}. \\
  \hline
  \texttt{*(}, \texttt{*)} & Символы \texttt{\{} и \texttt{\}}. \\
  \hline
  \texttt{*t} & Символ табуляции. \\
  \hline
  \texttt{**} & Символ \texttt{*}. \\
  \hline
  \texttt{*'}, \texttt{*"} & Одинарная и двойная кавычки. \\
  \hline
  \texttt{*n} & Перевод строки. \\
  \hline
\end{longtable}

\section{Требования к компилятору}

Разрабатываемый компилятор должен выполнять следующие функции:

\begin{enumerate}
  \item чтение исходного файла на языке B;
  \item удаление пробельных символов и комментариев;
  \item выделение токенов и проверка корректности литералов;
  \item построение таблицы символов;
  \item разбор глобальных определений, функций, локальных объявлений и операторов;
  \item построение абстрактного синтаксического дерева;
  \item проверка корректности операций над машинными словами и адресами;
  \item генерация ассемблерного кода x86-64 в синтаксисе GNU assembler;
  \item вызов ассемблера и компоновщика для получения исполняемого файла;
  \item поддержка режима вывода только ассемблерного файла и режима вывода объектного файла;
  \item выполнение набора регрессионных тестов.
\end{enumerate}

\section{Вывод}

В аналитической части рассмотрены назначение компилятора, особенности языка B, теоретические основы трансляции, основные лексические элементы и функциональные требования к транслятору. Грамматика поддерживаемого языка вынесена в приложение А. Реализация соответствует основным конструкциям, описанным в руководстве по B: внешним определениям, функциям, \texttt{auto}, \texttt{extrn}, выражениям, указательной арифметике, векторам, условным операторам, циклам, \texttt{switch}, \texttt{goto} и \texttt{return}.
